% Standalone source-record documentation; no research attempts.
\documentclass[11pt]{article}
\usepackage[margin=25mm]{geometry}
\usepackage{fontspec,unicode-math}
\setmainfont{Latin Modern Roman}
\setmathfont{Latin Modern Math}
\usepackage{amsmath,amssymb,mathtools,xurl,hyperref}
\hypersetup{colorlinks=true,urlcolor=blue,linkcolor=blue}
\setcounter{secnumdepth}{0}
\setlength{\parindent}{0pt}
\setlength{\parskip}{5pt}
\setlength{\emergencystretch}{3em}
\newcommand{\sref}[2]{\hyperlink{source:#1:#2}{[S#2]}}
\newcommand{\cataloguescope}[1]{\paragraph{Source record.}#1}
\newcommand{\floor}[1]{\left\lfloor#1\right\rfloor}
\begin{document}
\section{A semantics-bound certified bridge between Lean and computer algebra}\label{statement}

\subsection{Definitions and mathematical statement}

\paragraph{Shared notation.}Standard number systems are $\mathbb N,\mathbb Z,\mathbb Q,\mathbb R,\mathbb C$. Specialized notation and index ranges follow the cited source and the definitions below.


A computer algebra system performs symbolic computations on mathematical expressions and algebraic structures. A formal theorem prover checks derivations in an explicit logical framework. A certificate is additional data whose correctness can be checked; connecting a symbolic computation to a proof requires relating its encoded input and result to formal mathematical definitions.
\paragraph{Record scope and qualifications.}
The catalogue title is retained as metadata. Named certificates, contracts, bounds or frameworks in that title are not asserted as results of the original source extract.
The extract does not specify a computation, logic, interface, trust boundary, or verification theorem. It includes a software-design discussion as well as possible implementation work.

Catalogue screening limitation: Workshop planning or software infrastructure agenda; the supplied passage does not identify a determinate mathematical target.

The source snapshot labels this record partially solved; this is catalogue metadata rather than independent verification of a mathematical result or the current completion of a project.
\paragraph{Statement recorded in the supplied source.}
The group has the express goal of finding ways in which current computer algebra systems can interact with formal theorem provers, such as Lean, and vice versa. The guiding motivation for this project is to discuss to what extent Lean should operate as a CAS in it's own right. This groups concrete goals could be anything from the actual development of CAS technology in Lean, or could be to create an argument about the extent to which Lean should be used as a computer algebra system at all.
The primary focus of this group should be to decide to what extent these mathematical technologies should mix and to come up with ways to implement these solutions. \sref{statement}{2}
\paragraph{Source status.}
The supplied catalogue labels this source record partially solved; that metadata was checked against the supplied snapshot on 2026-10-11. The label is preserved for provenance and is not an independent verification of present open status or mathematical correctness.
\subsection{Short English statement}
Explore interactions between computer algebra systems and formal theorem provers, including Lean.
\subsection{Sources}
\begin{itemize}
\item[\hypertarget{source:statement:1}{S1}] Ulam AI and the UnsolvedMath Contributors, UnsolvedMath: A Curated Collection of Open Mathematics Problems, catalogue entry. Catalogue curation and metadata: CC BY 4.0. The source extract is retained; definitions, English summaries and qualifications are editorial additions.
\url{https://huggingface.co/datasets/ulamai/UnsolvedMath}.
\item[\hypertarget{source:statement:2}{S2}] AIM Workshop Problem Lists, source linked by the catalogue.
\url{http://aimpl.org/cyberinfrastructure/6/}.
\end{itemize}
\paragraph{Preparation and review.}
The original source and UnsolvedMath contributors retain attribution for the source extract. Codex (GPT-6; exact serving revision not exposed) assisted with definitions, English summaries, scope notes and formatting. No new human mathematical authorship or independently verified solution is asserted. Catalogue identity checked against the complete supplied registry and current repository. Distinct upstream records, including repeated questions and broad or damaged extracts, are retained for complete coverage of the source CS category. A repeated record is not a new independent conjecture. Complete mathematical-equivalence and current-literature review has not been performed. Preparation and final review follow contribution prompts C and D, with source gaps disclosed. Structural and rendering checks do not establish a new theorem or open status.
\label{end:statement}
\end{document}
